This paper measures whether a maximum-likelihood fit of a multivariate
exponential-kernel Hawkes process reports intervals that behave as a user is
entitled to assume, and it ships the answer as code. Over a pre-registered grid of
\HsixCells{} cells and \HsixReplications{} simulated replications per cell, the
inverse-Hessian standard errors understate sampling spread: the ratio of mean
reported standard error to empirical standard deviation, pooled over the twelve
two-type cells with at least 2000 events, is \HsixQ{} $\pm$ \HsixQSe{}, and the
four-type satellite agrees in direction (\HsixSatQ{} $\pm$ \HsixSatSe{}). Two
hypotheses were refuted: H6.3 predicted that coverage falls monotonically toward
criticality, and H6.4 predicted that timestamp quantisation biases the decay
estimate upward. The estimator is undefined
in \HsixUndefined{} of \HsixCells{} cells, where the observation window holds too
few decay times to identify the decay rate. The library therefore refuses to
report an interval there, and where the map measures a shortfall it switches to a
parametric bootstrap rather than reporting the interval it knows to be too narrow;
the switch is tested in a second pre-registered study and its result is reported
either way.
